\chapter{Known Issues}
\label{ch:known-issues}

\section{Purpose and method}

This chapter records defects, inconsistencies and behavioural surprises that are
\emph{present in the code as analysed}. It does not duplicate security findings,
which are catalogued in Chapter~\ref{ch:security}, nor the maintenance backlog in
Chapter~\ref{ch:technical-debt}. Each issue is given an identifier, a category, a
severity and an evidence pointer. Severity reflects impact on a user of the
product: \textbf{High} means wrong or broken generated output in a plausible
workflow, \textbf{Medium} means confusion or a manual workaround, \textbf{Low}
means cosmetic or documentation-only.

Some issues are documented by the authors themselves (the extension README's
``Known Limitations'' and the CHANGELOG), and these are marked accordingly.

\section{Issue register}

\begin{longtable}{p{1.2cm}p{3.3cm}p{1.6cm}p{1.1cm}p{6.4cm}}
\caption{Known issues}\\
\toprule
\textbf{ID} & \textbf{Category} & \textbf{Severity} & \textbf{Status} &
\textbf{Description and evidence} \\
\midrule

KI-01 & Generated styling & High & Open & shadcn-generated screens use design tokens
such as \texttt{bg-background}, \texttt{text-foreground},
\texttt{border-border} and \texttt{ring-ring}, but the repository ships no
Tailwind/NativeWind theme configuration defining those tokens. A generated shadcn
screen can therefore render unstyled until the user adds the token definitions.
Evidence: shadcn templates under \texttt{packages/registry/templates}; no
\texttt{tailwind.config.*} in the repository. \\
KI-02 & Missing dependency & High & Open & The \texttt{glass} variants of
\texttt{app/home}, \texttt{auth/login} and \texttt{onboarding/welcome} import
\texttt{expo-blur}, which no recipe declares in \texttt{dependencies}; the package
is not installed. Evidence: \texttt{*.glass.hbs} template imports; recipe JSON
dependency arrays. \\
KI-03 & Interactive generation & High & Open (documented) & Auth hooks generated by
\texttt{auth/login} and \texttt{auth/register} are placeholders: they wait a fixed
delay and then always report success. The shadcn login form only
\texttt{console.log}s the submitted data and does not call the hook at all.
Shipping without replacement yields a non-functional login. Evidence:
\texttt{hooks.use-login.hbs}, \texttt{hooks.use-register.hbs},
\texttt{components.login-form.shadcn.hbs}. \\
KI-04 & Advertised-but-unimplemented options & Medium & Open (documented) &
\texttt{tamagui}, \texttt{vector-icons} and \texttt{custom} are offered in the
extension settings and exist in the schema, but no recipe implements them.
Selecting one finds no matching recipe. Evidence: extension
\texttt{package.json} enums vs recipe \texttt{supported} arrays. \\
KI-05 & CLI version reporting & Medium & Open & The CLI's \texttt{--version} output
is a hard-coded literal rather than the package version, so it will drift from
\texttt{package.json} on release. Evidence:
\texttt{packages/cli/src/index.ts}. \\
KI-06 & Silent overwrite & Medium & Open & \texttt{add} overwrites any existing file
at a recipe path without confirmation or backup; only primitive files are
skip-aware. Evidence: \texttt{cli/src/utils/writeFiles.ts}. (Also logged as
SEC-16.) \\
KI-07 & Variant validation & Medium & Open & An unknown \texttt{--variant} value is
silently ignored and the default template is produced, with no warning. Evidence:
\texttt{core/src/generate.ts} \texttt{resolveTemplateString}. \\
KI-08 & Silent fallback & Medium & Open & When a configured remote registry is
unreachable, the extension silently uses its built-in catalog; the user cannot
tell that their configured registry was ignored. Evidence: extension
\texttt{registry.ts} fallback path. \\
KI-09 & Stale product document & Medium & Open & \texttt{documents/ProjectDescription.md}
understates the catalog and implies variants are unsupported, contradicting the
code (24 recipes; variants present). Evidence: document vs
\texttt{packages/registry/data/**} and \texttt{registry/index.json}. \\
KI-10 & Partial-failure consistency & Low & Open & If dependency installation
fails, generated files have already been written and the manifest may not be
updated, leaving files and manifest disagreeing. Evidence: ordering in
\texttt{cli/src/commands/add.ts}. (Also logged as SEC-17.) \\
KI-11 & Credential logging & Low & Open & Generated auth code logs submitted
credentials to the console. Evidence: auth hook templates. (Also logged as
SEC-12.) \\
KI-12 & Manifest tolerance & Low & Open & A malformed or unreadable
\texttt{.expo-snips.json} is treated as an empty manifest rather than reported, so
corruption is silent. Evidence: \texttt{cli/src/utils/manifest.ts}. \\
KI-13 & Generated artefacts retained in the tree & Low & Open & Two build
outputs are kept in the working tree: the built registry
(\texttt{registry/}, 25 files) and the offline bundle
(\texttt{packages/registry/bundle.json}) are \emph{tracked} by git, while four
packaged extensions (\texttt{expo-snippets-0.1.0..0.1.3.vsix}) sit untracked in
the extension directory (excluded by a \texttt{*.vsix} rule in
\texttt{.gitignore}). Tracking generated files invites drift from source.
Evidence: \texttt{git ls-files}, \texttt{.gitignore}. \\
KI-14 & Registry host unverifiable & Low & Unknown & The CLI's default registry
host \texttt{registry.expo-snips.dev} has no hosting configuration in the
repository, so its existence and ownership cannot be verified here. Evidence:
registry resolution constants; no CNAME or pages config. \\
KI-15 & Generated registry tracked in git & Low & Open & Because
\texttt{registry/} is tracked \emph{and} regenerated on every push to the default
branch, the deployed output and the committed copy can diverge if the deploy
workflow's build differs from a manual build. The single source of truth is
\texttt{packages/registry/data/}. Evidence: \texttt{git ls-files registry};
\texttt{deploy-registry.yml}. \\
\bottomrule
\end{longtable}

\section{Resolved since the initial documentation}

The 1.0 release of this document set (17 October 2026) did not record two
template defects that were discovered the following day while generating every
screen for the repository's screenshot set. Both were fixed in the source
templates on 19 October 2026 and are documented in \texttt{CHANGELOG.md}
(Unreleased section).

\begin{longtable}{p{1.6cm}p{4.6cm}p{4.4cm}p{3.4cm}}
\caption{Resolved template defects (19 October 2026)}\\
\toprule
\textbf{ID} & \textbf{Defect} & \textbf{Evidence} & \textbf{Resolution} \\
\midrule

RS-01 & JSX object literals written as \texttt{\{\{ uri: ... \}\}},
\texttt{initialRegion=\{\{...\}\}} and \texttt{style=\{\{...\}\}} were parsed by
Handlebars as helper calls, so \texttt{chat/conversations},
\texttt{content/article-reader} and \texttt{map/location} crashed at generation
(\emph{Missing helper: ``uri:''}; a parse error in the map template).
\textbf{Status:} IMPLEMENTED (fixed). &
\texttt{templates/chat/conversations/components.conversation-item.hbs},
\texttt{templates/content/article-reader/app.article.[id].hbs},
\texttt{templates/map/location/app.map.index.hbs}. &
Escaped literals (\texttt{\textbackslash\{\{}). All three recipes generate correctly;
\texttt{pnpm validate} reports 0 errors. \\
RS-02 & Component props and calendar helpers were emitted without TypeScript
types, so generated output failed \texttt{tsc --noEmit} in strict projects
(implicit \texttt{any}).
\textbf{Status:} IMPLEMENTED (fixed). &
\texttt{components.conversation-item.hbs}, \texttt{components.event-item.hbs},
\texttt{components.location-sheet.hbs}, \texttt{app.calendar.index.hbs}. &
Typed props/parameters added behind the catalog's existing
\texttt{\{\{#if typescript\}\}} guard. \\
\bottomrule
\end{longtable}

These entries are listed here rather than in the open issue register because
they no longer reproduce against the current tree.

\section{Author-acknowledged limitations}

The extension README's ``Known Limitations'' section explicitly states: no tests;
no CI/CD for the extension (packaged manually); limited icon/styling options
(\texttt{tamagui}, \texttt{vector-icons}, \texttt{custom} unsupported); offline
by default; and VS Code only (the extension and CLI are independently installed).
These align with KI-02 through KI-05 above and with the findings in
Chapter~\ref{ch:testing} and Chapter~\ref{ch:cicd}.

\textbf{Status:} IMPLEMENTED (documented by the authors). This is a positive
practice --- the limitations are stated rather than hidden.

\section{Corrections required to prior documentation}

Two specific corrections are needed to bring the pre-existing
\texttt{ProjectDescription.md} into agreement with the code:

\begin{longtable}{p{5.2cm}p{4.2cm}p{4.6cm}}
\caption{Documentation corrections}\\
\toprule
\textbf{Claim as written} & \textbf{Code reality} & \textbf{Correction} \\
\midrule

An earlier recipe count (stated as 19). & 24 recipes in
\texttt{packages/registry/data}. & Update the count to 24. \\
Variants described as unsupported. & \texttt{app/home}, \texttt{auth/login} and
\texttt{onboarding/welcome} define a \texttt{glass} variant. & Document variant
support and the three recipes that use it. \\
\bottomrule
\end{longtable}

\section{Issue density by area}

\begin{table}[H]
\centering
\caption{Issues by category}
\begin{tabular}{p{5cm}p{4cm}p{5cm}}
\toprule
\textbf{Category} & \textbf{Count} & \textbf{Highest severity} \\
\midrule
Generated output correctness & 3 (KI-01, KI-02, KI-03) & High \\
Configuration / UX & 4 (KI-04, KI-05, KI-07, KI-08) & Medium \\
Data integrity & 3 (KI-06, KI-10, KI-12) & Medium \\
Documentation / hygiene & 5 (KI-09, KI-11, KI-13, KI-14, KI-15) & Medium \\
\bottomrule
\end{tabular}
\end{table}

\section{Summary}

\textbf{Status:} the product is a functioning, coherent tool, but three
high-severity issues affect the \emph{output} it produces: shadcn styling tokens
without theme configuration (KI-01), an undeclared \texttt{expo-blur} import on
glass variants (KI-02), and placeholder authentication code that a developer
might ship unmodified (KI-03). These are the items to address first because they
undermine the product's core promise of ``production-quality screens''.
